This master's thesis addresses the design and implementation of a~mobile
application that enables operators of an~industrial production line to
quickly identify a~specific device, monitor it, and remotely control it
using a~unique visual identifier. An~analysis of available commercial
platforms showed that none of them simultaneously integrate a~native
mobile application, scanning of unique identifiers for pairing with
a~device, and a~context-aware intelligent assistant with a~custom
knowledge base. The thesis addresses these gaps by proposing the
Scan2Control system, in which a~mobile application implemented in Flutter
automatically connects the operator to a~specific line by scanning an
ArUco identifier and subsequently allows them to send control commands.
The backend layer, formed by a~pair of FastAPI
services, communicates with a~Siemens S7-1200 programmable logic
controller through Node-RED and an~MQTT broker, while the entire
infrastructure is containerized using Docker and deployed on an~edge
device in the production environment. The system includes the
intelligent assistant PLEXO, built on the Retrieval-Augmented Generation
architecture, which answers technical questions from operators based on
a~knowledge base stored in the ChromaDB vector database. Testing in
a~laboratory environment, and partially also in an~industrial setting,
confirmed the fulfilment of all functional, non-functional and
domain-specific requirements, including command latency significantly
below the defined threshold and reliable detection of identifiers across
a~practical range of distances, angles and lighting conditions.